\textbf{China Intelligent-connected Car Assessment Programme，}中文官方文件标题里叫：\textbf{《中国智能网联汽车技术规程（C-ICAP）》，是一个自愿/评价性质的第三方测评。一般来说，车企新车型上市，想证明智驾能力的时候会去测。}

C-NCAP / Euro NCAP 的测评主目标是“证明这辆车有多安全”，如果要“证明这辆车的智驾有多强”，就需要C-ICAP出马。

\begin{quote}
  \textbf{注意:} 举个例子假设你的车有城市 NOA。测 C-NCAP / Euro NCAP，可能重点看：AEB 能不能避免撞车、撞人车道保持是否安全驾驶员是否被有效监控系统失效时有没有安全兜底但它不一定完整证明：城市路口会不会博弈能不能识别复杂车道线能不能处理加塞、绕行、变道导航辅助体验是否连贯泊车和座舱体验是否智能这些更像 C-ICAP 的评价范围。
\end{quote}

其中

C-ICAP 的行车辅助分三层：

\begin{enumerate}
  \item \textbf{基础行车辅助}
  \item \textbf{高速领航辅助}
  \item \textbf{城市领航辅助}
\end{enumerate}

在 2024 年版修订版里，\textbf{基础行车辅助是核心必测项}；\textbf{高速领航辅助、城市领航辅助是审核项}，正式测评中企业可在高速领航和城市领航里二选一提供第三方测试报告。
